% =============================================================================
% Technical paper: Bayesian magnetic search, mutual information, dipole inverse
% Overleaf: upload algorithm_paper.zip, main.tex, pdfLaTeX
% =============================================================================
\documentclass[11pt,a4paper]{article}

\usepackage[margin=2.2cm]{geometry}
\usepackage{amsmath,amssymb,amsthm}
\usepackage{booktabs}
\usepackage{graphicx}
\usepackage{hyperref}
\usepackage{xcolor}
\usepackage{enumitem}
\usepackage{caption}
\usepackage{fancyhdr}
\usepackage{listings}
\usepackage{float}
\usepackage{setspace}

\hypersetup{colorlinks=true,linkcolor=blue!40!black,urlcolor=blue!50!black,citecolor=blue!40!black}
\pagestyle{fancy}
\fancyhf{}
\fancyhead[L]{\small Bayesian Magnetic Search --- Algorithm Paper}
\fancyhead[R]{\small Multi-ASV / ROS~2}
\fancyfoot[C]{\thepage}
\renewcommand{\headrulewidth}{0.35pt}
\setstretch{1.08}

\newtheorem{proposition}{Proposition}
\newtheorem{remark}{Remark}

\title{\textbf{Bayesian Grid Search, Mutual-Information Hunt,\\
and Inverse-Cube Dipole Localization}\\[0.4em]
\large An Engineering Algorithm Paper for the Cooperative ASV Magnetic Survey Stack}
\author{Internal study document matching \texttt{version\_3} implementation\\
(\texttt{bayes\_core.py}, \texttt{info\_gain.py}, \texttt{dipole\_fit.py})}
\date{\today}

\begin{document}
\maketitle

\begin{abstract}
This paper explains \emph{exactly} the three estimation/planning algorithms used
in the cooperative ASV magnetic survey: (1)~a discrete Bayesian occupancy
search on a lake grid with a soft \(1/r^{3}\) detection likelihood;
(2)~a mutual-information (MI) waypoint planner that chooses the next pose to
maximize expected entropy reduction; (3)~a continuum inverse-cube dipole fit
(Levenberg--Marquardt) that refines \((t_x,t_y,A)\) from strong anomaly samples.
We state the physical inverse-cube law, the probabilistic model, the update
algorithms, why they are consistent, the high-level parameters, and how a
\emph{second} ASV uses the same dipole/belief products to verify. The goal is
study material: after this paper one should be able to explain to others
\emph{what} we do, \emph{why} it works, and \emph{which knobs} matter.
\end{abstract}

\tableofcontents
\vspace{0.6em}

%==============================================================================
\section{Engineering picture (what the stack is doing)}
%==============================================================================

The magnetic target is unknown. Each ASV measures a scalar anomaly feature
\(a\) (nT) after filtering/calibration. We never give estimators the planted
truth \(\mathbf{t}^{\star}\). Three complementary estimators run in parallel:

\begin{enumerate}[leftmargin=*]
  \item \textbf{Bayesian grid (coarse, probabilistic).}
        Maintain a probability mass \(b_k\) that the target sits in cell \(k\).
        A HIT near a cell multiplies that cell by a high detection probability;
        cells far away are down-weighted. This \emph{concentrates} belief.
  \item \textbf{Mutual information (where to go next).}
        Given current \(b\), score candidate poses by how much they would
        \emph{shrink entropy in expectation}. Move there (hunt).
  \item \textbf{Dipole inverse (fine, physics).}
        Once many strong \(a_i\) exist, fit the continuum model
        \(a=A/(r^{3}+s^{3})\) for location \((t_x,t_y)\) and amplitude \(A\).
        This is the metre-class fix published on \texttt{/swarm/belief/fix}.
\end{enumerate}

Multi-ASV: all boats dump anomalies into \emph{one} shared map
(\texttt{bayes\_fusion}). The discoverer hunts/spirals; a \emph{peer} verifier
orbits the declared candidate using the same belief + dipole products.

\begin{lstlisting}
SEARCH (strips) -> Bayes HITs concentrate b
 -> MI hunt -> spiral densify samples
 -> declare candidate (centroid)
 -> peer VERIFY (same dipole/Bayes evidence)
 -> COMPLETE; score ||t_fix - t*||
\end{lstlisting}

%==============================================================================
\section{Physical model: inverse-cube (why \(1/r^{3}\))}
%==============================================================================

\subsection{From magnetostatics to our scalar plant}
A magnetic dipole \(\mathbf{m}\) at \(\mathbf{t}\) produces field (SI)
\[
\mathbf{B}(\mathbf{r})
=\frac{\mu_0}{4\pi}
\left(
\frac{3(\mathbf{m}\cdot\hat{\mathbf{r}})\hat{\mathbf{r}}-\mathbf{m}}{r^{3}}
+\frac{2}{3}\mathbf{m}\,\delta^{3}(\mathbf{r})
\right),
\quad \mathbf{r}=\mathbf{p}-\mathbf{t}.
\]
Away from the origin the leading term scales as \(\mathbf{B}\propto 1/r^{3}\).
For a \emph{vertical} dipole observed mainly in \(B_z\) near the surface, the
scalar anomaly is well approximated by a radial \(1/r^{3}\) envelope.

\subsection{Regularized plant used in code}
To keep the field finite overhead (\(r=0\)) we use the \emph{softened} law
implemented in \texttt{dipole.py} / \texttt{dipole\_fit.py}:
\begin{equation}
\label{eq:dipole}
a(\mathbf{p};\,\mathbf{t},A,s)
=\frac{A}{r^{3}+s^{3}},
\qquad r=\lVert\mathbf{p}-\mathbf{t}\rVert_{2}.
\end{equation}
Parameters (YAML): \(A=4\times10^{5}\,\mathrm{nT\cdot m^{3}}\),
\(s=20\,\mathrm{m}\) \(\Rightarrow\) peak overhead \(A/s^{3}=50\,\mathrm{nT}\).
At \(r=s=20\,\mathrm{m}\), \(a=A/(2s^{3})=25\,\mathrm{nT}\).

\begin{remark}[Inverse law]
This is \textbf{not} Coulomb \(1/r^{2}\). Magnetic dipole far-field is
\(\mathbf{B}\propto 1/r^{3}\). The \(s^{3}\) term is a numerical/physical
softening (finite sensor height / source size), not a change of the far-field
exponent.
\end{remark}

Noise model (plant): Earth background \(\sim 45{,}000\,\mathrm{nT}\) plus
\(\mathcal{N}(0,3^{2})\,\mathrm{nT}\). HIT threshold \(15\,\mathrm{nT}\) is
\(\sim 5\sigma\) above noise.

%==============================================================================
\section{Bayesian grid search}
\label{sec:bayes}
%==============================================================================

\subsection{State}
Discretize the \(300\times300\,\mathrm{m}\) lake into cells of size
\(\Delta=10\,\mathrm{m}\) (origin \((-150,-150)\)): \(N=30\times30=900\) cells.
Let \(T\in\{1,\ldots,N\}\) be the unknown target cell. Belief is a PMF
\[
b_k \triangleq \Pr(T=k),\qquad \sum_{k=1}^{N} b_k=1.
\]
\textbf{Prior:} uniform \(b_k^{(0)}=1/N\) (\texttt{BeliefMap} constructor).

\subsection{Observation model (HIT / MISS / ABSTAIN)}
From cleaned anomaly \(a_{\mathrm{cl}}\) (\texttt{classify\_observation}):
\[
z=
\begin{cases}
\mathrm{HIT} & a_{\mathrm{cl}}\ge \tau_{\mathrm{hit}}=15\,\mathrm{nT},\\
\mathrm{MISS} & a_{\mathrm{cl}}\le \tau_{\mathrm{miss}}=5\,\mathrm{nT},\\
\mathrm{ABSTAIN} & \text{otherwise (or uncalibrated)}.
\end{cases}
\]
With \texttt{hit\_only: true}, MISS is treated as ABSTAIN (no update). This is
conservative: weak ``no detection'' does not falsely erase mass.

\subsection{Likelihood inspired by \(1/r^{3}\)}
If the ASV is at \(\mathbf{x}\) and the hypothesized target is cell centre
\(\mathbf{c}_k\), \(d_k=\lVert\mathbf{x}-\mathbf{c}_k\rVert\). Detection
probability (\texttt{detection\_probability}):
\begin{equation}
\label{eq:pdet}
P_{\mathrm{det}}(d)
=P_{\mathrm{bg}}+(P_{\max}-P_{\mathrm{bg}})
\frac{d_{1/2}^{3}}{d^{3}+d_{1/2}^{3}},
\end{equation}
with \(P_{\mathrm{bg}}=0.05\), \(P_{\max}=0.95\), \(d_{1/2}=30\,\mathrm{m}\).
Interpretation:
\begin{itemize}[leftmargin=*]
  \item Far away: \(P_{\mathrm{det}}\to P_{\mathrm{bg}}\) (false-alarm floor).
  \item Overhead: \(P_{\mathrm{det}}\to P_{\max}\) (almost sure detect).
  \item Same \emph{shape} as the dipole envelope \(d^{3}/(d^{3}+d_{1/2}^{3})\).
\end{itemize}
We set
\[
\Pr(z=\mathrm{HIT}\mid T=k)=P_{\mathrm{det}}(d_k),\qquad
\Pr(z=\mathrm{MISS}\mid T=k)=1-P_{\mathrm{det}}(d_k).
\]

\subsection{Bayes update (the algorithm)}
One calibrated HIT at pose \(\mathbf{x}\):
\begin{equation}
\label{eq:bayes-hit}
b_k^{+}
\;\propto\;
b_k^{-}\, P_{\mathrm{det}}(\lVert\mathbf{x}-\mathbf{c}_k\rVert),
\qquad
b^{+}=b^{+}/\textstyle\sum_{j} b_{j}^{+}.
\end{equation}
A MISS (if enabled) multiplies by \(1-P_{\mathrm{det}}\). This is exactly
\texttt{BeliefMap.update} in \texttt{bayes\_core.py}.

\begin{proposition}[Why HITs concentrate on the true cell]
Assume the ASV records a HIT at \(\mathbf{x}\) near the true target cell
\(k^{\star}\). Then \(P_{\mathrm{det}}(d_{k^{\star}})\gg P_{\mathrm{det}}(d_{\ell})\)
for distant cells \(\ell\). After renormalization,
\(b_{k^{\star}}\) increases relative to far cells. Repeated independent HITs
around \(k^{\star}\) drive \(b\to\) a unimodal blob (MAP = peak cell).
\end{proposition}

\begin{remark}[Recursive Bayesian filter]
This is a discrete Bayes filter / histogram filter on a static target.
It is the same recursion as occupancy-grid ``target search'' (Stone / Bayesian
search theory): likelihood \(\times\) prior, then normalize. It does
\emph{not} assume Gaussianity.
\end{remark}

\subsection{Summaries published to ROS}
\begin{itemize}[leftmargin=*]
  \item \textbf{Peak} \(\arg\max_k b_k\) \(\to\) \texttt{/swarm/belief/peak},
        \(p^{\star}=b_{\mathrm{peak}}\).
  \item \textbf{Centroid} of cells with \(b_k\ge 0.5\,b_{\max}\), mass-weighted:
        \[
        \hat{\mathbf{t}}_{\mathrm{c}}
        =\frac{\sum_{k\in\mathcal{H}} b_k\mathbf{c}_k}{\sum_{k\in\mathcal{H}} b_k},
        \quad
        \sigma^{2}
        =\frac{\sum_{k\in\mathcal{H}} b_k\lVert\mathbf{c}_k-\hat{\mathbf{t}}_{\mathrm{c}}\rVert^{2}}
              {\sum_{k\in\mathcal{H}} b_k}.
        \]
        Sub-cell estimate; \(\sigma\) is published as
        \texttt{centroid\_spread}.
\end{itemize}
Mission uses \textbf{centroid preferred over peak} as the hunt/declare estimate.

\subsection{Why this works (engineering argument)}
\begin{enumerate}[leftmargin=*]
  \item The likelihood \eqref{eq:pdet} is \emph{matched} to the physics
        \eqref{eq:dipole}: both decay as \(1/d^{3}\). Matched filters beat
        generic isotropic kernels.
  \item Uniform prior + HIT-only updates cannot invent mass far from real HITs;
        false HITs are suppressed by the \(15\,\mathrm{nT}\approx5\sigma\) gate.
  \item Multiple ASVs share one \(b\) (independent observations, same \(T\)):
        information adds. That is why ASV2/ASV3 lawnmower HITs help ASV1's map.
  \item Limitation: grid quantization (\(\Delta=10\,\mathrm{m}\)) \(\Rightarrow\)
        peak jumps between cell centres. Centroid + dipole LS fix that.
\end{enumerate}

\subsection{Pseudocode}
\begin{lstlisting}
b <- uniform(N cells)
for each MagAnomaly (x,y,a_cl) from any ASV:
  if not calibrated or 5 < a_cl < 15: continue   # ABSTAIN
  if a_cl >= 15:   # HIT
    for each cell k:
      b[k] *= P_det(||(x,y)-c_k||)
    b <- b / sum(b)
publish peak, p*, centroid, spread
\end{lstlisting}

%==============================================================================
\section{Mutual-information hunt}
\label{sec:mi}
%==============================================================================

\subsection{Information-theoretic objective}
Shannon entropy of the belief
\[
H(b)=-\sum_{k} b_k\log b_k
\]
(\texttt{belief\_entropy}; natural log). A good next pose \(\mathbf{c}\) is one
that, \emph{in expectation}, reduces \(H\) the most. That expected reduction
\emph{is} mutual information between target cell \(T\) and the next binary
observation \(Z\in\{\mathrm{HIT},\mathrm{MISS}\}\):
\begin{equation}
\label{eq:mi}
I(T;Z\mid\mathbf{c})
=H(T)-H(T\mid Z,\mathbf{c})
=H(b)-\mathbb{E}_{Z}\!\left[H(b'\mid Z,\mathbf{c})\right].
\end{equation}

\subsection{Predictive HIT probability}
At candidate pose \(\mathbf{c}\),
\[
p_{\mathrm{HIT}}(\mathbf{c})
=\sum_{k} b_k\, P_{\mathrm{det}}(\lVert\mathbf{c}-\mathbf{c}_k\rVert)
\]
(\texttt{expected\_posterior\_entropy}: \texttt{p\_hit = sum(b*p\_det)}).
Then
\[
b^{\mathrm{HIT}}_k \propto b_k P_{\mathrm{det},k}(\mathbf{c}),\qquad
b^{\mathrm{MISS}}_k \propto b_k\bigl(1-P_{\mathrm{det},k}(\mathbf{c})\bigr),
\]
and
\[
\mathbb{E}[H(b')]
= p_{\mathrm{HIT}} H(b^{\mathrm{HIT}})
+(1-p_{\mathrm{HIT}}) H(b^{\mathrm{MISS}}).
\]
\texttt{mutual\_information\_gain} returns \(H(b)-\mathbb{E}[H(b')]\).

\begin{proposition}[MI is non-negative]
For any \(\mathbf{c}\), \(I(T;Z\mid\mathbf{c})\ge 0\), with equality iff \(Z\)
is independent of \(T\) (e.g.\ \(P_{\mathrm{det}}\) constant in \(k\)). Hence
the planner never ``expects to lose information.''
\end{proposition}

\subsection{Candidate generation (what we actually search)}
We do \emph{not} optimize over the whole lake (too expensive). Ring candidates
around current pose (\texttt{info\_gain\_radii}: \(10,20,30\,\mathrm{m}\),
\(16\) angles) clipped to coverage bounds. Planner returns
\(\arg\max_{\mathbf{c}} I(T;Z\mid\mathbf{c})\), replanned every \(5\,\mathrm{s}\).

If the belief grid is not ready, fallback: pick the ring point closest to the
current peak (with a small move penalty).

\subsection{Why MI hunt works}
\begin{itemize}[leftmargin=*]
  \item Lawn mower is \emph{open-loop coverage}; MI is \emph{closed-loop}:
        go where the next reading most distinguishes competing cells.
  \item Near a unimodal blob, MI typically pulls the boat \emph{toward the
        uncertain rim / closer to the peak}, not randomly.
  \item Same \(P_{\mathrm{det}}\) as Bayes \(\Rightarrow\) planner and filter
        are \textbf{model-consistent} (no planner/filter mismatch).
  \item Exit to spiral when MI collapses, near peak (\(\le25\,\mathrm{m}\)),
        step budget \(45\), or \(p^{\star}\ge0.5\) near estimate.
\end{itemize}

\subsection{Pseudocode}
\begin{lstlisting}
candidates = ring(current, radii={10,20,30}, 16 angles) inside bounds
for c in candidates:
  I[c] = H(b) - (pHIT*H(b_HIT) + pMISS*H(b_MISS))
go to argmax I
\end{lstlisting}

%==============================================================================
\section{Dipole inverse estimation (Levenberg--Marquardt)}
\label{sec:dipole}
%==============================================================================

\subsection{Problem}
Bayes gives a \emph{cell}. Physics gives a \emph{continuum} location. Buffer
samples \(\{(x_i,y_i,a_i)\}\) with \(a_i\ge 10\,\mathrm{nT}\) (min 12, max 400)
and solve
\begin{equation}
\label{eq:ls}
\min_{t_x,t_y,A}
\sum_{i=1}^{n}
\left(
\frac{A}{r_i^{3}+s^{3}}-a_i
\right)^{2},
\quad r_i=\sqrt{(x_i-t_x)^{2}+(y_i-t_y)^{2}},\ s=20.
\end{equation}
This is nonlinear least squares. Code: \texttt{DipoleFitter.fit}.

\subsection{Jacobians (what LM actually uses)}
Let \(f_i=A/(r_i^{3}+s^{3})\). Then
\[
\frac{\partial f_i}{\partial A}=\frac{1}{r_i^{3}+s^{3}},
\qquad
\frac{\partial f_i}{\partial r}=-3A r_i^{2}/(r_i^{3}+s^{3})^{2}.
\]
Chain rule \(\partial r/\partial t_x=-d_x/r\) yields
\[
\frac{\partial f_i}{\partial t_x}
=\frac{3A r_i\, d_x}{(r_i^{3}+s^{3})^{2}},
\quad
\frac{\partial f_i}{\partial t_y}
=\frac{3A r_i\, d_y}{(r_i^{3}+s^{3})^{2}}
\]
(\(r>0\); at \(r=0\) spatial grads vanish, soft term keeps \(f\) finite).
LM solves \((J^{\top}J+\lambda D)\delta=-J^{\top}\mathbf{r}\) with residual
\(\mathbf{r}=\mathbf{f}-\mathbf{a}\), damping \(\lambda\), step capped at
\(40\,\mathrm{m}\). Success if RMS \(<20\,\mathrm{nT}\) and \(A>0\).

\subsection{Warm start}
If a guess \(\hat{\mathbf{t}}_{\mathrm{c}}\) (centroid) is available, use it;
else anomaly-weighted sample centroid. Strength guess \(A\approx a(r^{3}+s^{3})\)
from the strongest sample, or YAML \(4\times10^{5}\).

\subsection{Why dipole LS works (and when)}
\begin{itemize}[leftmargin=*]
  \item \textbf{Identifiability:} \(1/r^{3}\) curvature is unique enough that
        angular diversity (spiral/orbit) plus amplitude \(A\) can be recovered.
        A single radial pass is poorly conditioned --- that is why we spiral.
  \item \textbf{Matched model:} plant and fitter use the \emph{same}
        \(A/(r^{3}+s^{3})\) family \(\Rightarrow\) zero-bias in sim when
        samples surround \(\mathbf{t}^{\star}\) (observed best \(e_{\mathrm{fix}}\approx0\)).
  \item \textbf{Bayes vs dipole:} Bayes is robust to outliers and multi-ASV
        fusion; dipole is precise but needs strong, geometrically diverse
        samples. We use Bayes to \emph{get there}, dipole to \emph{score metres}.
\end{itemize}

%==============================================================================
\section{How the other ASV uses dipole / belief (peer verify)}
%==============================================================================

\begin{enumerate}[leftmargin=*]
  \item Discoverer finishes spiral, declares \textbf{centroid}
        \((c_x,c_y,p)\) on \texttt{/swarm/verify/declare}.
  \item \texttt{verify\_coordinator} assigns a \emph{different} ASV
        (\texttt{prefer\_other\_verifier: true}).
  \item Verifier flies an orbit about the candidate (radius \(\le20\,\mathrm{m}\),
        inland-capped) approaching from the opposite side.
  \item Each anomaly sample counts as a confirmation iff
        \emph{all} hold (\texttt{verify\_core}):
        at-site \(\le30\,\mathrm{m}\), peak within \(50\,\mathrm{m}\) of
        candidate, \(a_{\mathrm{cl}}\ge15\,\mathrm{nT}\), \(p^{\star}\ge0.30\).
  \item Need 4 confirms \(\to\) COMPLETE. Meanwhile dipole LS keeps updating
        from \emph{all} ASV samples (shared buffer in \texttt{bayes\_fusion}).
\end{enumerate}
So the peer does \textbf{not} run a different physics law. It re-samples the
same inverse-cube field and the same Bayes map from a new geometry --- that is
the statistical independence that makes verify meaningful.

%==============================================================================
\section{Parameter catalog (how it is modelled)}
%==============================================================================

\begin{center}
\small
\begin{tabular}{@{}llp{7.2cm}@{}}
\toprule
\textbf{Symbol / YAML} & \textbf{Value} & \textbf{Meaning} \\
\midrule
\(\Delta\), \texttt{cell\_size\_m} & \(10\,\mathrm{m}\) & Bayes cell \\
map size / origin & \(300\,\mathrm{m}\), \((-150,-150)\) & Lake grid \\
\(P_{\mathrm{bg}},P_{\max}\) & \(0.05,0.95\) & FA floor / max \(P_{\mathrm{det}}\) \\
\(d_{1/2}\) & \(30\,\mathrm{m}\) & Range where \(P_{\mathrm{det}}\) is mid \\
\(\tau_{\mathrm{hit}},\tau_{\mathrm{miss}}\) & \(15,5\,\mathrm{nT}\) & HIT/MISS gates \\
\texttt{hit\_only} & true & Ignore MISS updates \\
centroid frac & \(0.5\) & High-prob blob cutoff \\
\(A,s\) plant+fit & \(4\times10^{5}\), \(20\,\mathrm{m}\) & Dipole amplitude/soft \\
dipole min \(a\), \(n\) & \(10\,\mathrm{nT}\), \(12\) & LS buffer gate \\
MI radii / angles & \(10,20,30\,\mathrm{m}\), \(16\) & Hunt candidates \\
MI replan & \(5\,\mathrm{s}\) & Hunt rate \\
enter TARGET \(p^{\star}\) & \(0.25\times4\) ticks & Leave lawnmower \\
spiral max / inset & \(80\), \(10\,\mathrm{m}\) & Inland densify \\
verify \(n\) / orbit & \(4\), \(20\,\mathrm{m}\) & Peer confirm \\
\bottomrule
\end{tabular}
\end{center}

\paragraph{How to retune (engineering).}
Raise \(\tau_{\mathrm{hit}}\) \(\to\) fewer false HITs, slower detect.
Raise \(d_{1/2}\) \(\to\) longer-range belief influence (more smeared blob).
Smaller \(\Delta\) \(\to\) finer peak, more CPU. Larger MI radii \(\to\) bolder
hunt steps. Dipole \(s\) must match plant or LS is biased.

%==============================================================================
\section{End-to-end consistency argument}
%==============================================================================

\begin{enumerate}[leftmargin=*]
  \item \textbf{Physics \(\to\) likelihood:} \(B\propto 1/r^{3}\) \(\Rightarrow\)
        \(P_{\mathrm{det}}\) uses the same exponent.
  \item \textbf{Likelihood \(\to\) Bayes:} HIT update is Bayes' rule on a
        categorical \(T\); \(b\) remains a valid PMF.
  \item \textbf{Bayes \(\to\) MI:} MI uses the \emph{same} \(P_{\mathrm{det}}\)
        to predict \(b'\); planner is Bayes-optimal for one-step entropy.
  \item \textbf{Samples \(\to\) dipole:} LS inverts the \emph{same} \(A/(r^{3}+s^{3})\)
        used to generate \(a\) in sim \(\Rightarrow\) consistent estimator.
  \item \textbf{Peer ASV:} independent pose, same field, same \(b\) and LS
        buffer \(\Rightarrow\) geometric diversity + confirmation.
\end{enumerate}
This chain is why we can tell others: ``it is not three unrelated hacks; it is
one inverse-cube world model used as plant, likelihood, planner, and fitter.''

%==============================================================================
\section{Study questions (for explaining to others)}
%==============================================================================
\begin{enumerate}[leftmargin=*]
  \item Why \(1/r^{3}\) and not \(1/r^{2}\)? (dipole vs monopole)
  \item What does \(s\) do at \(r=0\)? (finite peak \(A/s^{3}\))
  \item Write one HIT Bayes update for two cells. Which cell grows?
  \item Why HIT-only? What would aggressive MISS updates do far from the target?
  \item Show \(I(T;Z)\ge 0\). When is \(I=0\)?
  \item Why rings of \(10/20/30\,\mathrm{m}\) instead of full-lake MI?
  \item Why spiral before dipole LS? (angular diversity / conditioning)
  \item Difference between peak cell, centroid, and dipole fix?
  \item Why can ASV1 verify a target ASV3 discovered? (shared \(b\), shared field)
  \item If \(\tau_{\mathrm{hit}}\) were \(3\,\mathrm{nT}\) (noise level), what
        fails? (false HITs, wrong blob, early false COMPLETE)
\end{enumerate}

%==============================================================================
\section{Code map}
%==============================================================================
\begin{center}
\begin{tabular}{@{}ll@{}}
\toprule
Algorithm & File \\
\midrule
\(P_{\mathrm{det}}\), Bayes HIT/MISS, centroid & \texttt{boat\_mapping/bayes\_core.py} \\
Fusion node + dipole buffer & \texttt{boat\_mapping/bayes\_fusion.py} \\
LM dipole fit + Jacobians & \texttt{boat\_mapping/dipole\_fit.py} \\
Plant dipole & \texttt{boat\_sensing/dipole.py} \\
MI planner & \texttt{boat\_mission/info\_gain.py} \\
FSM / declare / verify & \texttt{mission\_manager.py}, \texttt{verify\_core.py} \\
Parameters & \texttt{mapping.yaml}, \texttt{mission.yaml}, \texttt{sensing.yaml} \\
\bottomrule
\end{tabular}
\end{center}

\begin{thebibliography}{9}
\bibitem{shannon} C.~E.~Shannon, ``A mathematical theory of communication,''
\emph{Bell Syst.\ Tech.\ J.}, 1948.
\bibitem{lm} K.~Levenberg, ``A method for the solution of certain non-linear
problems in least squares,'' \emph{Q.\ Appl.\ Math.}, 1944.
\bibitem{stone} L.~D.~Stone, \emph{Theory of Optimal Search}, 1975.
\bibitem{fossen} T.~I.~Fossen, \emph{Handbook of Marine Craft Hydrodynamics and
Motion Control}, Wiley, 2011.
\bibitem{ros2} Open Robotics, ROS~2 Jazzy documentation.
\end{thebibliography}

\end{document}
